\begin{exercise}{38}
Let $F\colon\mathcal C\to\mathcal D$ be left-adjoint to $G\colon\mathcal D\to\mathcal C$.
Show that if $G$ is lax monoidal, then $F$ canonically refines to an oplax monoidal functor.
Vice versa, show that if $F$ is oplax monoidal, then $G$ canonically refines to a lax monoidal functor.

\begin{proof}
Let $\eta,\varepsilon$ be the unit and counit. For maps
$a\colon U\to GV$ and $b\colon FU\to V$, write
\[
 a^\flat=\varepsilon_VF(a),\qquad
 b^\sharp=G(b)\eta_U.
\]
These are inverse natural bijections by Exercise~\ref{ex:9}.
Use $X,Y\in\mathcal C$ and $A,B\in\mathcal D$ as variables.
Write $\mu,\mu_0$ for lax structure maps of $G$, and
$\delta,\delta_0$ for oplax structure maps of $F$.
The two canonical transfers are
\begin{equation}\label{eq:oplax-mate-formula}
 \delta_{X,Y}=
  \bigl(\mu_{FX,FY}(\eta_X\otimes\eta_Y)\bigr)^\flat,
 \qquad \delta_0=(\mu_0)^\flat,
\end{equation}
and, in the other direction,
\begin{equation}\label{eq:lax-mate-formula}
 \mu_{A,B}=
  \bigl((\varepsilon_A\otimes\varepsilon_B)\delta_{GA,GB}\bigr)^\sharp,
 \qquad \mu_0=(\delta_0)^\sharp.
\end{equation}
The binary maps are natural by binaturality of transposition.
The induced structures satisfy the identities
\begin{equation}\label{eq:monoidal-mates}
 \begin{aligned}
  G(\delta_{X,Y})\eta_{X\otimes Y}
   &=\mu_{FX,FY}(\eta_X\otimes\eta_Y),\\
  \varepsilon_{A\otimes B}F(\mu_{A,B})
   &=(\varepsilon_A\otimes\varepsilon_B)\delta_{GA,GB},\\
  G(\delta_0)\eta_{\bbone_{\mathcal C}}
   &=\mu_0,\qquad
    \varepsilon_{\bbone_{\mathcal D}}F(\mu_0)=\delta_0.
 \end{aligned}
\end{equation}
One identity in each direction is the defining mate identity.
For example, when $\mu$ is constructed from $\delta$,
\[
 \bigl(\mu_{FX,FY}(\eta_X\otimes\eta_Y)\bigr)^\flat
 =(\varepsilon_{FX}\otimes\varepsilon_{FY})
   \delta_{GFX,GFY}F(\eta_X\otimes\eta_Y)
 =\delta_{X,Y},
\]
by naturality of $\delta$ and
$\varepsilon_{FX}F(\eta_X)=\mathrm{id}_{FX}$.
When $\delta$ is constructed from $\mu$, naturality of $\mu$
and the other triangle identity give
\[
 \bigl((\varepsilon_A\otimes\varepsilon_B)\delta_{GA,GB}\bigr)^\sharp
 =\mu_{A,B}\bigl(G(\varepsilon_A)\eta_{GA}
                \otimes G(\varepsilon_B)\eta_{GB}\bigr)
 =\mu_{A,B}.
\]
The unit identities follow by the same triangle identities.

We verify coherence for the transfer from oplax $F$ to lax $G$.
Write $\lambda,\rho$ for the inward unit isomorphisms, the
inverses of those in Definition~1.2.30.
For $A,B,C\in\mathcal D$, let $L,R$ be the two composites
in the lax associativity diagram for $G$, with source
$(GA\otimes GB)\otimes GC$ and target $G(A\otimes(B\otimes C))$.
Put $E=\varepsilon_A\otimes(\varepsilon_B\otimes\varepsilon_C)$.
Using \eqref{eq:monoidal-mates} and naturality of
$\delta,\varepsilon$, their adjuncts are
\[
 \begin{aligned}
  L^\flat&=E\,\alpha^{\mathcal D}_{FGA,FGB,FGC}
   (\delta_{GA,GB}\otimes\mathrm{id})
   \delta_{GA\otimes GB,GC},\\
  R^\flat&=E\,(\mathrm{id}\otimes\delta_{GB,GC})
   \delta_{GA,GB\otimes GC}
   F(\alpha^{\mathcal C}_{GA,GB,GC}).
 \end{aligned}
\]
These are equal by oplax associativity of $\delta$.
Injectivity of transposition proves lax associativity.
Similarly, the adjuncts of the two unit diagrams reduce to
\[
 \begin{aligned}
  \lambda_A^{\mathcal D}(\delta_0\otimes\varepsilon_A)
   \delta_{\bbone_{\mathcal C},GA}
   &=\varepsilon_A F(\lambda_{GA}^{\mathcal C}),\\
  \rho_A^{\mathcal D}(\varepsilon_A\otimes\delta_0)
   \delta_{GA,\bbone_{\mathcal C}}
   &=\varepsilon_A F(\rho_{GA}^{\mathcal C}).
 \end{aligned}
\]
These equalities follow from the oplax unit axioms and
naturality of the unit isomorphisms, proving both lax unit axioms.

For the other transfer, apply this verified construction to
$G^{\mathrm{op}}\colon\mathcal D^{\mathrm{op}}\to\mathcal C^{\mathrm{op}}$
left-adjoint to $F^{\mathrm{op}}$, with unit
$\varepsilon^{\mathrm{op}}$ and counit $\eta^{\mathrm{op}}$.
The opposite monoidal categories keep the tensor factors in
the same order and have associator $(\alpha^{-1})^{\mathrm{op}}$
and inward unitors $(\lambda^{-1})^{\mathrm{op}}$,
$(\rho^{-1})^{\mathrm{op}}$.
A lax structure on $G$ becomes an oplax structure on
$G^{\mathrm{op}}$. The construction gives a lax structure on
$F^{\mathrm{op}}$, hence an oplax structure on $F$.
Reversing the arrow compositions gives exactly
\eqref{eq:oplax-mate-formula} and transports all the checked
associativity and unit diagrams to the required diagrams.
\end{proof}
\end{exercise}
